\section{Packets, localization, and fractal estimates}
\label{sec:preliminaries}

We first fix the packet conventions and the analytic estimates used in the
transverse argument.  There are two packet classes: compactly supported smooth
atoms, which are convenient for the later geometric arguments, and exact
Schr\"odinger packets, which permit the use of restriction estimates.  The latter
have rapidly decreasing spatial tails.  Keeping these two classes distinct will
also make the localization errors explicit.

\subsection{Conventions and the compact atom class}

For frequency data $v\in L^2(\R^2)$ supported in a fixed bounded set, put
\begin{equation}
 Ev(x,t)=\int_{\R^2}e^{i(x\cdot\xi-t|\xi|^2)}v(\xi)\,d\xi.
 \label{eq:prelim-extension}
\end{equation}
Thus extension input norms are frequency-space $L^2$ norms; the Fourier
normalization in Section~\ref{sec:introduction} differs only by the fixed factor
$(2\pi)^{-2}$.  In particular, with that Fourier convention,
$\|f\|_2^2=(2\pi)^{-2}\|\widehat f\|_2^2$.
All frequency sets below lie in fixed bounded boxes unless another scale is
specified.  The direction associated with $\xi$ is the line parallel to
$(2\xi,1)$, or its unit representative with positive last coordinate.  On a
fixed bounded frequency set, distances between frequencies and between these
unit representatives are comparable.

Unless stated otherwise, all measures are nonnegative Borel measures on space-time $\R^3$.  A ball
bound is required for balls with arbitrary centers.  Its normalizing parameter
may be assumed positive, since the zero case is immediate.  Lattice partitions
are half-open; a supremum on a cell may be taken on its closure.  Fixed
enlargements, boundedly overlapping grids, and fixed changes of frequency boxes
cost constants independent of the large scales and of the number of packets.
All distribution thresholds are positive.

An estimate with a free factor $S^\epsilon$ is asserted for each fixed
$\epsilon>0$, with a constant that may depend on $\epsilon$.  We can therefore
use a smaller exponent to absorb a fixed number of logarithms.  On a
countersequence, $o(1)$ denotes a quantity tending to zero; simultaneous
statements at countably many fixed parameters are obtained by subsequences and
diagonal choices.  Parameters reused in different arguments have only their
local meaning.

\begin{definition}[Compact atoms]
\label{def:prelim-atoms}
Let $L\geq1$, and set $R=L^2$.  An atom with label $\alpha$ is the field
\begin{equation}
 A_\alpha(x,t)=L^{-1}
 e^{i(x\cdot\omega_\alpha-t|\omega_\alpha|^2)}
 H_\alpha\left(\frac{x-b_\alpha-2\omega_\alpha t}{L},
                  \frac{t}{L^2}\right).
 \label{eq:src5}
\end{equation}
The frequencies $\omega_\alpha$ lie in a fixed square.  The labels obey a
joint phase-space bound: every unit ball in $\R^2\times\R^2$ contains at most
a fixed number of the pairs
\[
 \left(\frac{b_\alpha}{L},L\omega_\alpha\right).
\]
The profiles $H_\alpha(y,\tau)$ are supported in a fixed bounded box, and
have uniform bounds for sufficiently many derivatives.  More precisely, each
application requires a finite integer $J$ and fixed bounds
$\|\partial^\beta H_\alpha\|_\infty\leq C_\beta$ for $|\beta|\leq J$.
The integer $J$ is chosen after the fixed parameters of that application; it
does not grow with $L$ or with the array cardinality.  A packet field is a finite
sum $\sum_\alpha c_\alpha A_\alpha$, and its coefficient energy is
$\sum_\alpha|c_\alpha|^2$, without weights.
\end{definition}

The joint phase-space condition does not require bounded multiplicity of
frequencies alone.  In particular, arbitrarily many spatial translates of one
frequency are allowed.  Constants in packet estimates may depend on the fixed
support, multiplicity, and derivative bounds in
Definition~\ref{def:prelim-atoms}.  Whenever a later estimate permits an
arbitrary fixed decay order, sufficiently many of these finitely many
derivatives are to be fixed first.

Suppose that $\mu$ is restricted to $|t|\leq C L^2$ and satisfies
\begin{equation}
 \mu(B(z,r))\leq m r^2 \qquad(z\in\R^3,\ r\geq1).
 \label{eq:prelim-ball-bound}
\end{equation}
The additional sparsity assumption, with the fixed choice $M=100$, is
\begin{equation}
 \int_{\R^3}
 \left(1+\frac{|x-b_\alpha-2\omega_\alpha t|}{L}\right)^{-M}
 \,d\mu(x,t)
 \leq\frac{mL^3}{D}
 \quad\text{for every label used.}
 \label{eq:src6}
\end{equation}
Other sufficiently large fixed values of $M$ would work.  The ball bound alone
implies \eqref{eq:src6} with $D=1$ up to a fixed constant.  Indeed, the
$\rho$-neighborhood of one trajectory on this time slab is covered by
$O(1+L^2/\rho)$ balls of radius $C\rho$.  For $\rho=2^kL$, its mass is at most
\[
 Cm\bigl(2^kL^3+2^{2k}L^2\bigr).
\]
Summing these bounds against $2^{-Mk}$, $k\geq0$, proves the assertion.

The normalization in \eqref{eq:src5} is chosen to be stable under cap
rescaling.  For labels in a square of side $h^{-1}$ centered at $v$, where
$1\leq h\leq L$, set
\begin{equation}
 (x',t')=\left(\frac{x-2vt}{h},\frac{t}{h^2}\right),\qquad
 L'=\frac Lh,\quad b'_\alpha=\frac{b_\alpha}{h},\quad
 \omega'_\alpha=h(\omega_\alpha-v).
 \label{eq:prelim-atom-rescaling}
\end{equation}
An exact substitution gives
\begin{equation}
 A_\alpha(x,t)=h^{-1}e^{i(x\cdot v-t|v|^2)}A'_\alpha(x',t'),
 \label{eq:prelim-atom-covariance}
\end{equation}
where $A'_\alpha$ has the same profile and the parameters in
\eqref{eq:prelim-atom-rescaling}.  The coefficients are unchanged, and
\[
 \left(\frac{b'_\alpha}{L'},L'\omega'_\alpha\right)
 =\left(\frac{b_\alpha}{L},L\omega_\alpha-Lv\right),
\]
so the phase-space ball counts are unchanged as well.  If $\mu'$ is the
pushforward of $\mu$, an inverse image of an $r$-ball is covered by $O(h)$
balls of radius $Chr$.  Hence $\mu'$ has ball parameter $m'=Cmh^3$.
The integrand in \eqref{eq:src6} is invariant under this change, and
$m'(L')^3=CmL^3$, so its sparsity parameter is preserved up to the allowed
fixed constants.  On a time slab of length $O((L')^2)$ centered at zero, a
compact packet meets only $O(1)$ spatial squares of side $(L')^2$.  Thus the
energies assigned to these squares sum with bounded overlap.

\subsection{Exact packets and spatial localization}

The corresponding energy statement for exact fields requires controlling the
tails rather than declaring the packets compactly supported.

\begin{lemma}[Exact packet localization]
\label{lem:packets}
Let $s\geq1$, let $v$ have frequency support in a fixed bounded set, and
consider $|t|\leq Cs^2$.  There is a decomposition
$v=\sum_{\omega,n}a_{\omega,n}v_{\omega,n}$, with
$\omega\in s^{-1}\mathbb{Z}^2$ in a fixed bounded enlargement and
$b_n=sn$, $n\in\mathbb{Z}^2$, such that
\begin{align}
 \sum_{\omega,n}|a_{\omega,n}|^2&\leq C\|v\|_2^2,
 \label{eq:prelim-packet-energy}\\
 \left\|\sum_{(\omega,n)\in\mathcal I}
                  a_{\omega,n}v_{\omega,n}\right\|_2^2
 &\leq C\sum_{(\omega,n)\in\mathcal I}|a_{\omega,n}|^2
 \quad\text{for every subfamily }\mathcal I,
 \label{eq:prelim-subseries-energy}\\
 |Ev_{\omega,n}(x,t)|&\leq C_Ns^{-1}
 \left(1+\frac{|x-b_n-2t\omega|}{s}\right)^{-N}
 \quad(N\geq1).
 \label{eq:prelim-packet-decay}
\end{align}
The constants are independent of the number of terms retained.  The
phase-space label counts are bounded, and every packet has frequency support
in an $O(s^{-1})$ neighborhood of $\omega$.

Fix $0<\eta<1$.  On a set $\Omega$ in this time slab, retain the packets
whose trajectories come within $s^{1+\eta}$ of $\Omega$ at the same time.
For every $A>0$, the sum of the omitted packets is bounded pointwise on
$\Omega$ by $C_{A,\eta}s^{-A}\|v\|_2$.
For a partition into spatial squares of side $s^2$, extended over the common
time slab, the retained data $v_Q$ can consequently be chosen so that
\begin{equation}
 \sum_Q\|v_Q\|_2^2\leq C\|v\|_2^2,
 \qquad
 |Ev-Ev_Q|\leq C_{A,\eta}s^{-A}\|v\|_2
 \quad\text{on }Q\times[-Cs^2,Cs^2].
 \label{eq:prelim-box-localization}
\end{equation}
The same conclusions hold on translated time slabs after modulating the
initial datum and writing the trajectories in the recentered time $t-t_0$.
\end{lemma}

\begin{proof}
Take a smooth bounded-overlap partition $\{\chi_\omega\}$ of frequency
space at resolution $s^{-1}$.  We can arrange that
$\chi_\omega(\omega+\eta/s)$ is supported in $[-1,1]^2$.  Put
\[
 g_\omega(\eta)=s^{-1}(\chi_\omega v)(\omega+\eta/s).
\]
Then $\|g_\omega\|_2=\|\chi_\omega v\|_2$.  On
$[-\pi,\pi]^2$, expand $g_\omega$ in the orthonormal Fourier basis:
\[
 g_\omega(\eta)=\sum_{n\in\mathbb{Z}^2}
 a_{\omega,n}(2\pi)^{-1}e^{-in\cdot\eta}.
\]
Choose one smooth cutoff $\psi$ supported in $(-\pi,\pi)^2$ and equal to
one on $[-1,1]^2$.  Multiplying the series by $\psi$ leaves $g_\omega$
unchanged, and gives the packet data
\[
 v_{\omega,n}(\xi)=\frac{s}{2\pi}
 \psi\bigl(s(\xi-\omega)\bigr)
 e^{-isn\cdot(\xi-\omega)}.
\]
Parseval, followed by the bounded overlap of the frequency partition, gives
\eqref{eq:prelim-packet-energy}.  Multiplication by $\psi$ is bounded on
$L^2$, and its translated frequency supports have bounded overlap.  Applying
these two facts to any selected Fourier coefficients proves
\eqref{eq:prelim-subseries-energy}.  This also proves convergence of the
decomposition in $L^2$.  Since all frequency supports remain in a fixed bounded
set, the corresponding extensions converge uniformly in $(x,t)$.

Changing variables $\xi=\omega+\eta/s$ gives
\begin{equation}
 Ev_{\omega,n}(x,t)=\frac{s^{-1}}{2\pi}
 e^{i(x\cdot\omega-t|\omega|^2)}
 \int \psi(\eta)
 e^{i((x-sn-2t\omega)/s)\cdot\eta-i(t/s^2)|\eta|^2}\,d\eta.
 \label{eq:prelim-exact-profile}
\end{equation}
For $|t|/s^2\leq C$, all derivatives of
$\psi(\eta)e^{-i(t/s^2)|\eta|^2}$ are bounded.  Repeated integration by
parts proves \eqref{eq:prelim-packet-decay}.  In particular the profile in
\eqref{eq:prelim-exact-profile} is Schwartz in its first argument, uniformly
on the prescribed time interval.

For clarity, the tail bound is uniform even for infinite packet expansions.
At a fixed point, there are $O(s^2)$ frequency centers.  For each one the
normalized spatial labels form a translate of $\mathbb{Z}^2$, and, for $N>1$,
\[
 \sum_{n:\ |x-sn-2t\omega|>s^{1+\eta}}
 \left(1+\frac{|x-sn-2t\omega|}{s}\right)^{-2N}
 \leq C_Ns^{-\eta(2N-2)}.
\]
The factor $s^{-2}$ from the squared packet bound cancels the
$O(s^2)$ frequency count.  Cauchy--Schwarz and
\eqref{eq:prelim-packet-energy} therefore bound the discarded field by
$C_Ns^{-\eta(N-1)}\|v\|_2$.  Choose $N$ sufficiently large in terms of
$A$ and $\eta$.

Finally, a bounded-velocity trajectory traverses distance $O(s^2)$ in the
time slab.  Since $s^{1+\eta}\leq s^2$, its padded trajectory meets only
$O(1)$ squares of side $s^2$.  Each coefficient is consequently assigned to
only $O(1)$ of the subfamilies defining $v_Q$.
Equation~\eqref{eq:prelim-subseries-energy} proves their summed energy
bound, and the preceding tail estimate proves the second part of
\eqref{eq:prelim-box-localization}.  To center the time interval at $t_0$,
replace $v(\xi)$ by $e^{-it_0|\xi|^2}v(\xi)$ before decomposing it.
\end{proof}

In applications involving polynomially many boxes or a measure of polynomial
total mass, the exponent $A$ is chosen large enough that all these tail errors
remain smaller than any prescribed negative power.  If only a subpower
padding is desired on a countersequence, first fix a positive padding exponent
for each required estimate and then let it decrease through a diagonal choice.
This does not assert a uniform tail constant as $\eta\downarrow0$.

\subsection{The refined fractal estimate and unit-cube suprema}

We use the following specialization of the refined fractal estimate of
Du--Zhang.  We include the reductions from its uniform occupancy formulation,
since both the weights and the suprema will matter later.

\begin{lemma}[Refined fractal estimate]
\label{lem:fractal}
Let $S\geq1$.  Let $\mathcal Q$ be a collection of distinct unit lattice
cubes in a ball of radius $CS$ in $\R^3$.  Suppose that
\begin{equation}
 \#\{Q\in\mathcal Q:Q\cap B(z,r)\ne\varnothing\}
 \leq\gamma r^2\quad(1\leq r\leq CS),
 \qquad
 \#\{Q\in\mathcal Q:Q\cap B(z,\sqrt S)\ne\varnothing\}
 \leq\lambda_0
 \label{eq:prelim-fractal-counts}
\end{equation}
for all centers $z$, with fixed enlargements of the radius range allowed.
If $v$ has frequency support in a fixed bounded set, then
\begin{align}
 \int_{\bigcup\mathcal Q}|Ev|^2\,dx\,dt
 &\leq C_\epsilon S^\epsilon
       (\gamma\lambda_0S)^{1/3}\|v\|_2^2,
 \label{eq:prelim-fractal-integral}\\
 \sum_{Q\in\mathcal Q}\sup_Q|Ev|^2
 &\leq C_\epsilon S^\epsilon
       (\gamma\lambda_0S)^{1/3}\|v\|_2^2.
 \label{eq:prelim-fractal-sup}
\end{align}
These estimates are uniform under translations of the collection of cubes.
\end{lemma}

\begin{proof}
For unit-ball frequency support, Theorem~1.6 of \cite{DZ19}, specialized to
spatial dimension $2$ and counting exponent $2$, gives
\[
 \|Ev\|_{L^2(X)}
 \leq C_\epsilon S^\epsilon
       \gamma^{1/6}\lambda^{1/6}S^{1/6}\|v\|_2
\]
when every occupied square-root-scale lattice cube contains a comparable
number $\lambda$ of the unit cubes of $X$.  Its time sign is changed by
reflection, and its physical-space input norm is converted to the
frequency-space norm by Plancherel.  A bounded frequency enlargement is
reduced to unit frequency by a fixed parabolic dilation and bounded cube
coverings.  None of these operations changes a large-scale exponent.

For the asserted upper occupancy version, take an ambient radius comparable
to $S$ whose square root is an integer, large enough to contain the given
ball.  Its square-root grid then contains whole unit lattice cubes.  The
occupancy of each parent is at most $C\lambda_0$, since such a parent is
covered by a fixed number of $\sqrt S$-balls.  Sort the parents by their
dyadic occupancy $\lambda$.  Every resulting union retains the counting
bound $C\gamma r^2$.  Apply the quoted theorem to each union, square it,
and sum.  The geometric sum of $\lambda^{1/3}$ is at most
$C\lambda_0^{1/3}$; using a logarithmic bound would also suffice after
decreasing the chosen small exponent.  Renaming that exponent proves
\eqref{eq:prelim-fractal-integral}.  Ball centers can be translated at the
cost of replacing $v(\xi)$ by
$e^{i(y\cdot\xi-u|\xi|^2)}v(\xi)$, which preserves its norm and support.

To obtain suprema, let $F=Ev$.  Its space-time Fourier transform is supported
on a fixed compact portion of the paraboloid.  Choose a Schwartz function
$\rho$ whose Fourier transform is one on that compact set, so that
$F=F*\rho$.  For every unit cube $Q$ and every $z\in Q$,
\begin{equation}
 \sup_{y\in Q}|F(y)|
 \leq C_N\int_{\R^3}(1+|w|)^{-N}|F(z+w)|\,dw.
 \label{eq:prelim-reproducing}
\end{equation}
Indeed, taking the supremum over $y-z$ in a fixed bounded set in the
convolution formula leaves a rapidly decreasing majorant for the kernel.
Cauchy--Schwarz with this integrable weight, followed by integration over
$z\in Q$ and summation in $Q$, gives
\[
 \sum_{Q\in\mathcal Q}\sup_Q|F|^2
 \leq C_N\int_{\R^3}(1+|w|)^{-N}
       \int_{\bigcup\mathcal Q}|Ev(z+w)|^2\,dz\,dw.
\]
For each fixed $w=(y,u)$ the inner field is the extension of the modulated
datum above.  Apply \eqref{eq:prelim-fractal-integral} on the unchanged
domain and integrate the weight.  This proves
\eqref{eq:prelim-fractal-sup} without enlarging the ambient scale according
to the translation.
\end{proof}

The same reproducing argument applies separately to factors in a product.
For example, if $F_i=Ev_i$, $i=1,2,3$, then
\begin{equation}
 \sum_{Q\in\mathcal Q}\prod_{i=1}^3\sup_Q|F_i|
 \leq C_N\int_{(\R^3)^3}\prod_{i=1}^3(1+|w_i|)^{-N}
 \int_{\bigcup\mathcal Q}\prod_{i=1}^3|Ev_i(z+w_i)|\,dz
 \,dw_1\,dw_2\,dw_3.
 \label{eq:prelim-product-suprema}
\end{equation}
To see this, multiply \eqref{eq:prelim-reproducing} for the three fields,
integrate over each unit cube, and use Tonelli.  For each fixed triple of
shifts, the three exact fields are tested at the same integration point on
the same domain.  Their data have independent unimodular modulations, which
preserve supports, norms, and frequency transversality.  This is the version
used when applying multilinear restriction to products of cube suprema.

\subsection{Weights and cap rescaling}

\begin{lemma}[Weighted refined estimate]
\label{lem:prelim-weighted}
Let $\nu$ be supported in a ball of radius $CS$, $S\geq1$, and suppose
\begin{equation}
 \nu(B(z,r))\leq mr^2\quad(r\geq1),\qquad
 \nu(B(z,\sqrt S))\leq mSd
 \quad(z\in\supp\nu),
 \label{eq:prelim-weighted-hypotheses}
\end{equation}
where $m,d>0$.  For bounded frequency support,
\begin{equation}
 \int |Ev|^2\,d\nu
 \leq C_\epsilon S^\epsilon mS^{2/3}d^{1/3}\|v\|_2^2.
 \label{eq:prelim-weighted-fractal}
\end{equation}
The same bound holds for
$\sum_Q\nu(Q)\sup_Q|Ev|^2$, summed over unit cubes.  The ball bound alone
gives this conclusion with $d=1$, up to a fixed constant.
\end{lemma}

\begin{proof}
Fixed-radius enlargements of the second bound in
\eqref{eq:prelim-weighted-hypotheses} follow by a bounded covering with
balls of radius $\sqrt S$ centered on the support.  This also extends it
to arbitrary centers, up to a fixed constant.  Every unit cube has mass
at most $Cm$.  Split the positive-mass cubes into classes
\[
 m\chi<\nu(Q)\leq2m\chi,
 \qquad \chi\in\{C_02^{-k}:k\geq0\},
\]
with harmless endpoint adjustments and fixed $C_0$.  Cubes in one class
have counting parameter $C\chi^{-1}$: the cubes meeting an $r$-ball have
disjoint masses inside a fixed enlargement of that ball.  Their
square-root-scale occupancy is at most $CSd\chi^{-1}$ by the same
argument.  Lemma~\ref{lem:fractal} therefore gives
\begin{align*}
 \sum_{Q\text{ in the }\chi\text{ class}}\nu(Q)\sup_Q|Ev|^2
 &\leq C_\epsilon m\chi S^\epsilon
   \bigl(\chi^{-1}\cdot Sd\chi^{-1}\cdot S\bigr)^{1/3}
   \|v\|_2^2\\
 &=C_\epsilon S^\epsilon mS^{2/3}d^{1/3}
   \chi^{1/3}\|v\|_2^2.
\end{align*}
The series $\sum_\chi\chi^{1/3}$ converges.  Summing proves both
assertions, including measures with arbitrarily small positive cube masses;
no lower mass cutoff is needed here.  Finally the first ball bound, evaluated
at $r=\sqrt S$, supplies $d=1$.
\end{proof}

\begin{lemma}[Weighted cap estimate]
\label{lem:cap}
Fix $\alpha_0<1/2$.  Suppose that $S\geq1$, that $\nu$ is supported in
a ball of radius $CS$ and satisfies $\nu(B(z,r))\leq mr^2$ for $r\geq1$,
and that $v$ is supported in a frequency ball $B(v_0,q)$, with $v_0$ in a
fixed bounded set and
\[
 S^{-\alpha_0}\leq q\leq1.
\]
Then, for every $\epsilon>0$,
\begin{equation}
 \int |Ev|^2\,d\nu
 \leq C_{\epsilon,\alpha_0}S^\epsilon
       mS^{2/3}q^{1/3}\|v\|_2^2.
 \label{eq:prelim-cap-estimate}
\end{equation}
The estimate also holds with the integral replaced by
$\sum_Q\nu(Q)\sup_Q|Ev|^2$ over unit cubes.  Bounded translations, and
averages of translations with a fixed integrable Schwartz weight, obey the
same bounds.
\end{lemma}

\begin{proof}
Translate the space-time support ball to the origin, absorbing the
translation in a unimodular factor on $v$.  Write
\[
 g(\eta)=q\,v(v_0+q\eta),\qquad
 (y,\tau)=\bigl(q(x-2v_0t),q^2t\bigr),\qquad S'=Sq^2.
\]
Then $\|g\|_2=\|v\|_2$, $\supp g\subset B(0,1)$, and
\begin{equation}
 Ev(x,t)=q\,e^{i(x\cdot v_0-t|v_0|^2)}Eg(y,\tau).
 \label{eq:prelim-cap-identity}
\end{equation}
Let $\nu'$ be the pushforward of $\nu$.  The inverse map is
\[
 t=q^{-2}\tau,\qquad x=q^{-1}y+2v_0q^{-2}\tau.
\]
Consequently, the inverse image of a ball of radius $r\geq1$ is a bounded
shear of a box with side lengths $r/q,r/q,r/q^2$.  It is covered by
$O(q^{-1})$ balls of radius $Cr/q$.  The original ball bound applies at
these radii and gives
\begin{equation}
 \nu'(B(z,r))\leq Cm q^{-3}r^2\qquad(r\geq1).
 \label{eq:prelim-rescaled-density}
\end{equation}
There is no Jacobian multiplier on this pushforward measure.

The rescaled time interval has length $O(S')$, but the spatial projection
has diameter $O(Sq)=O(S'/q)$.  Partition the latter into squares $Q$ of
side $S'$.  On the common time slab, use Lemma~\ref{lem:packets} at packet
width $s=\sqrt{S'}$, with any fixed padding exponent in $(0,1)$.  Let
$g_Q$ denote its localized data.  Then
\[
 \sum_Q\|g_Q\|_2^2\leq C\|g\|_2^2,
 \qquad |Eg-Eg_Q|\leq C_A(S')^{-A}\|g\|_2
 \quad\text{on the corresponding prism},
\]
where $A$ can be made arbitrarily large.  Each restricted measure lies in a
ball of radius $CS'$, so Lemma~\ref{lem:prelim-weighted} with $d=1$ and
\eqref{eq:prelim-rescaled-density} gives
\[
 \sum_Q\int_{Q\times[-CS',CS']}|Eg_Q|^2\,d\nu'
 \leq C_\epsilon(S')^\epsilon m q^{-3}(S')^{2/3}\|g\|_2^2.
\]
The sum uses the localized energies, rather than the number of spatial
squares.  This distinction is essential for the cap power.

For completeness, $\nu'(\R^3)=\nu(\R^3)\leq CmS^2$, whereas
$S'\geq S^{1-2\alpha_0}$.  Thus the integrated squared error, after
multiplication by the $q^2$ in \eqref{eq:prelim-cap-identity}, is at most
\[
 C_A m q^2 S^2(S')^{-2A}\|g\|_2^2.
\]
Choose $A$ in terms of $\alpha_0$ so that this is bounded by the desired
right side; bounded $S$ is absorbed in the constant.  For the principal term,
\begin{equation}
 q^2\,m q^{-3}(Sq^2)^{2/3}
 =mS^{2/3}q^{1/3}.
 \label{eq:prelim-cap-power}
\end{equation}
Since $(S')^\epsilon\leq S^\epsilon$, this proves
\eqref{eq:prelim-cap-estimate}.

Finally, square \eqref{eq:prelim-reproducing} by weighted
Cauchy--Schwarz, integrate in $z\in Q$ against $d\nu(z)$, and sum the
original unit cubes.  This bounds the sum of weighted suprema by
\[
 C_N\int_{\R^3}(1+|w|)^{-N}
           \int_{\R^3}|Ev(z+w)|^2\,d\nu(z)\,dw.
\]
For every shift the field is an exact extension of modulated data with the
same cap support and norm.  Apply the already proved cap estimate and
integrate the rapidly decreasing weight.  This proves the supremum
assertion and the asserted translated versions.
\end{proof}

\begin{remark}[Enlarged cubes in the inductive formulation]
\label{rem:prelim-enlarged-cubes}
The proof of Lemma~\ref{lem:fractal} is sufficient for all later applications.  For comparison,
we record an alternative reduction of \eqref{eq:prelim-fractal-integral} to
the inductive Proposition~3.1 of \cite{DZ19}.  Choose $K=S^\delta$ with $0<\delta<1/4$ sufficiently small
relative to the desired loss, and enlarge occupied unit cubes to occupied
$K^2$-cubes.  Selecting one original unit cube in each occupied large cube
shows, by ball covering, that the large cubes have counting parameter
$O(\gamma)$ for $r\geq K^2$, and square-root-scale occupancy
$O(\lambda_0)$.  Normalize $\|v\|_2=1$ and discard large cubes with
$L^6$ norm below $S^{-A}$.  There are polynomially many cubes, so their
total $L^2$ contribution is arbitrarily small when $A$ is sufficiently
large.  The remaining $L^6$ norms have only $O(\log S)$ dyadic classes,
since $|Ev|\leq C$ and each large cube has volume $K^6$.

After fixing a norm class, recompute its parent occupancies and sort them
dyadically.  On each resulting union $Y$ of $M_Y$ large cubes,
Proposition~3.1 gives
\[
 \|Ev\|_{L^6(Y)}
 \leq C_\epsilon M_Y^{-1/3}
       (\gamma\lambda_0S)^{1/6}S^\epsilon\|v\|_2.
\]
H\"older contributes
$|Y|^{1/3}=M_Y^{1/3}K^2$ on passing to $L^2$, cancelling $M_Y^{-1/3}$.
After squaring, the additional loss is $K^4$; the finitely many logarithmic
losses and this factor are absorbed by the choice of $\delta$ and a smaller
initial $\epsilon$.

The grids in this deduction can be made compatible as follows.  Increase
the ambient scale by a bounded factor so that
$S^{1/2}/S^{2\delta}$ is an integer.  Such scales have successive ratios
tending to one.  Use nested grids at these two scales and a bounded family
of translated parent grids.  For sufficiently large $S$, each original
unit cube can be assigned to a parent grid in which it is farther than
$2\sqrt3 K^2$ from the parent boundary; keep only the occupied large cubes
within that parent.  Split the original collection according to these
finitely many assignments.  The same original ball counts control each
part, while translation is absorbed into the datum as in
Lemma~\ref{lem:fractal}.  Bounded scales are covered by a fixed constant.
Thus neither boundary cubes nor rounding changes the resulting estimate.
\end{remark}
